\section*{How to Get Help}%
\label{sec:help}
We have lots of sources of help to facilitate your learning in CSSE 220. 
Don't resign yourself to struggling through this course alone. 
If you find that you have worked on something for 20 minutes 
without making any progress, 
it's probably time to seek help! 
Software development is a team sport, and 
the best developers know that a fresh set of eyes 
can often spot a problem right away. 

\subsection*{Instructor Email and Drop-In Hours:} 
See \href{https://docs.google.com/spreadsheets/d/1ip7XzOZyyDw10I6bsKDlkzAO0PowF4jPwrGiODMnQn8/edit?usp=sharing}{Ian's weekly schedule} (Section 01) or 
\href{https://docs.google.com/spreadsheets/d/13pt51VxLBnc-UkcqvvYSouB_YVTu_SEh5wCXk4I1-8c/edit?usp=sharing}{Aaron's weekly schedule} (Section 04)  
for student drop-in hours or to book an appointment. 
\textbf{The most effective and efficient way to get help} 
is to come to drop-in hours, and 
it's one of our favorite parts of the job, 
so please stop by. 
If you cannot attend drop-in hours, 
please book an appointment (see link in schedule) or email us at 
\href{mailto:luddenig@rose-hulman.edu}{luddenig@rose-hulman.edu} or 
\href{mailto:wilkin@rose-hulman.edu}{wilkin@rose-hulman.edu}. 
% (Ian: I check email twice daily Monday through Friday, 
% so you can expect a response within one business day.) 

% \subsection*{Teaching Assistant (TA) Help Hours:}
% Each section has an undergraduate teaching assistant (TA) 
% available to help you get unstuck during class. 
% For TA evening help hours, see the schedule in Moodle. 

\subsection*{Learning Center Tutors:}
The Learning Center in the Logan Library offers 
free peer tutoring for CSSE 220 and many other courses. 
See the Learning Center's 
\href{https://rosehulman.sharepoint.com/sites/LC}{My Rose-Hulman page} 
for this quarter's tutoring schedule and more info. 
They also offer \href{https://rosehulman.sharepoint.com/sites/LC/SitePages/SRT-Tutoring.aspx}{Sophomore Resident Tutors (SRTs)} in the Percopo classroom in the evenings. 
